\section{The first higher scalar Hall map}
\label{sch:sec:degree-four}

The scalar map for a plane and a line first has two incoming
contributions in degree $-4$. Let $s=(0,0,0)\in Z_3$, let
$G=\operatorname{Gr}(2,3)=\mathbb P^2$, and put
\[
 D_s=\left(Rb_{2,1*}a_{2,1}^*
                (\mathcal P_2\boxtimes\mathcal P_1)\right)_s.
\]
Stalks use ordinary atlas pullback. Let $\mathcal K$ be the pulled-back
coefficient family on $G$, so proper base change gives
$D_s=R\Gamma(G,\mathcal K)$.

\begin{proposition}[The incoming degree and its proper comparison]
\label{sch:prop:degree-four}
There is a canonical filtration of $H^{-4}(D_s)$ with successive
one-dimensional quotients
\[
 H^2(G,\mathcal H^{-6}\mathcal K),\qquad
 H^0(G,\mathcal H^{-4}\mathcal K).
\]
In particular $\dim H^{-4}(D_s)=2$.
Define smooth Milnor fibres
\[
 F=f_3^{-1}(1)\subset V_3,\qquad
 I=\{(A,B,C,H)\in E_{2,1}:f_3(A,B,C)=1\},
 \qquad q_1:I\longrightarrow F.
\]
The scalar coefficient map in degree $-4$ is identified with the
following composite:
\begin{equation}
 H^{-4}(D_s)\xrightarrow{\sim}H^5(I;\mathbb Q)
 \xrightarrow{\ q_{1*}\ }H^{13}(F;\mathbb Q)
 \xrightarrow{\sim}H^{\mathrm{BM}}_{22}(\mathcal C_3;\mathbb Q).
 \label{sch:eq:degree-four-gysin}
\end{equation}
Here $q_{1*}$ is the proper cohomological Gysin map with complex
orientations. The right isomorphism is the scalar support comparison
of Proposition~\ref{sch:prop:scalar}.
\end{proposition}

\begin{proof}
The rank-two scalar coefficient has degrees $-4,-2,1$, each of
rank one. The line coefficient has degree $-2$. Hence the cohomology
sheaves of $\mathcal K$ have degrees $-6,-4,-1$ and rank one.
They are locally constant: local frames trivialize the tautological
plane and quotient, and changes of frames preserve the scalar point.
Every local system on $\mathbb P^2$ is constant, since
$\mathbb P^2$ is simply connected.

In the hypercohomology spectral sequence
$E_2^{p,q}=H^p(G,\mathcal H^q\mathcal K)$, the terms of total
degree $-4$ are $(p,q)=(2,-6)$ and $(0,-4)$.
The first has no outgoing differential because no coefficient occurs
below $-6$. Its only possible incoming differential has source
$(0,-5)$, which is zero. The second has no incoming differential.
Its outgoing differential of length two has coefficient degree $-5$.
The length-three target is $(3,-6)$, where $H^3(\mathbb P^2)=0$.
Longer targets have coefficient degree below $-6$.
Thus both terms survive, giving the filtration and dimension.
The filtration does not select a splitting of the two contributions.

Write $E=E_{2,1}$ and $g=f_3|_E$.
The variety $E$ is a vector bundle of rank $21$ over $G$, and
$\dim_{\mathbb C}E=23$. Its block expression is $g=f_2+0$.
At potential value one the differential of the rank-two potential
is nonzero, so $I$ is smooth of complex dimension $22$.
Likewise $F$ is smooth of complex dimension $26$.
The restriction $q_1$ is proper, as a base change of the proper
incidence projection.

Positive scaling of the matrix entries contracts $E$ to its zero
section $G$ while preserving the sign of $\operatorname{Re}g$.
The same radial pair argument as in
Proposition~\ref{sch:prop:scalar} computes the scalar source by
\[
 D_s\simeq
 \operatorname{Fib}\bigl(R\Gamma(E,\mathbb Q)
                  \longrightarrow R\Gamma(I,\mathbb Q)\bigr)[10].
\]
The equality uses the rank-one zero-potential comparison, so it
retains the actual incoming vanishing-cycle coefficient.
Since $E$ retracts to $\mathbb P^2$, its cohomology in degrees
five and six vanishes. The relative-cochain exact sequence identifies
$H^{-4}(D_s)$ with $H^5(I)$.
The target is
$\operatorname{Fib}(\mathbb Q\to R\Gamma(F,\mathbb Q))[18]$.
Its degree $-4$ group is $H^{13}(F)$.

The ambient incidence Gysin map has degree eight. Restriction to
the smooth potential fibre is the proper Gysin map $q_{1*}$,
also of degree eight, since $26-22=4$.
The normal Thom construction and flag integration commute with
restriction of the potential. Therefore they induce a morphism
between the two relative-cochain sequences just used.
Both boundary identifications use the same normalized fibre convention.
The shift eight is even, so the square has no additional sign.
This proves the first two arrows of \eqref{sch:eq:degree-four-gysin}.
For $n=3$ and $j=-4$, the Borel--Moore index in
\eqref{sch:eq:scalar-bm} is $18-(-4)=22$, proving the last arrow.
\end{proof}

The projection of $I$ to its preserving pair $(A,C,H)$ has affine
fibres over the locus where the diagonal rank-two pair does not
commute. Forgetting $H$ from this pair locus is not proper.
Fix $H=\langle e_1,e_2\rangle$ and set
\[
 A=\operatorname{diag}(0,1,2),\qquad
 C_t=t e_{12}+e_{13},\qquad B_t=t^{-1}e_{21}\quad(t\ne0).
\]
These matrices preserve $H$ and satisfy
\[
 [C_t,A]=t e_{12}+2e_{13},\qquad
 \operatorname{Tr}(B_t[C_t,A])=1.
\]
The pair $(A,C_t)$ has a limit in $U_3$ at $t=0$, since
$[C_0,A]=2e_{13}\ne0$. Its two diagonal blocks commute at that
limit. Thus the limit lies outside the projected incidence locus.
The flag remains fixed, so no other limiting flag repairs this
missing point. The matrix $B_t$ diverges, consistently with the
properness of $q_1$ before this affine projection.

The image of the two-dimensional group in
\eqref{sch:eq:degree-four-gysin} therefore requires the proper
Milnor-fibre Gysin map, or a proved replacement retaining its support
data. The filtration specifies the two incoming contributions.
Their images and the rank of $q_{1*}$ in this degree remain to be
determined.
